\section{Peeling mosaics and conditional observation kernels}\label{sec:kernels}

This section proves conditional locality for limits of numerical observations on the actual discrete
flags and primal edges. The construction keeps the exact integer gluing information before taking
limits. Its inputs are the critical cutting and compatible-route results of
Section~\ref{sec:cutting} and the finite local order comparisons of Section~\ref{sec:walls}.
The main issue is that conditioning and weak convergence do not commute in general.
We generate the retained pieces independently in reference experiments and compare the
exact finite-map law to their product. The comparison density has a limit depending
only on macroscopic data that we subsequently condition on. Microscopic passage
observations can therefore retain arbitrary subsequential laws.

There are three steps. First, local slot identifications are encoded by integer
parameters; their image lattice splits across separated patches, with residues
retained. Second, the reference density comparison gives product kernels, including
when the full exploration tree and total area are specified. Third, the wall
comparisons provide one finite set of local chart choices for all the exterior
samples and all the countably many observation units under consideration.

\subsection{The primary input and buffered observations}

\begin{definition}[Primary tree input]\label{kern:primary-input}

Continue to use length unit \(P\) and step unit \(m_P\) of \eqref{eq:original-12}. Begin in pair law with
continuous-scale perimeter mixing (central first side scaled perimeter with a smooth positive scale density,
common typical phase); changes by the Palm weights described above will use density comparison. The signed
orientations are kept. Here is the primary tree input \({\cal T}\). In each canonical caller \(C\) use a
fresh length-sampled active source, trace toward the whole ordered boundary (without a priority order of
targets), making full branch jumps, continuing to inherited targets from the separating pinch. In a positive
disk, or in a negative at which the inherited original target interval becomes empty, start the rule afresh.
Keep this tree of length paths and boundary coordinates with side/order data, paired circumference orders,
and volume labels of domains; ambient conformal locations are not part of the abstract input. An overall
cyclic labeling on the central ring may be kept.

On the lattice use the all-boundary-slot version with exact
lists; branches before separation are shared, and the cut polygon pursued on switching is still fresh.
Exceptional absorptions for individual target slots can leave a microscopic terminal part for that path;
other paths continue by their own targets. Sample new sources independently in the as-yet-unexamined
polygons at the indicated births. Equivalently they can be assigned upon creation in whatever compatible
all-slot order is used. All fully reached macroscopic restarts will be away from such terminal exceptions.

\end{definition}

\begin{lemma}[Countable encoding and realization]\label{kern:encoding}
The primary input \(\mathcal T\) has a countable encoding by ordinary boundary targets, with
uniformly shrinking cells. Compatible lattice encodings converge jointly with volume labels and
admit the spatial realization of Section~\ref{sec:cutting}. Once their fractions are specified, paths to additional ordinary targets sampled from continuous
boundary densities are read from this same input.
\end{lemma}
\begin{proof}

For countable encoding it suffices in a continuum caller, and in joint lattice comparisons, to test a dense
ordinary sequence of target fractions (e.g. a randomized uniform sequence, or fixed ordinary fractions when
admissible). Use the first target still eligible each time on entering a negative with nonempty inherited
interval, and restart lists in newly born complete callers. Finite such routes, tested on compact survival
through matched pieces, converge jointly by \eqref{eq:original-12} and the compatible-route comparison of Proposition~\ref{cut:marked-targets}. At negative cuts along them an empty limiting inherited interval means a canonical return stabilized by Proposition~\ref{wall:returns}: its closure is disjoint even from the original caller rim. Thus it also has empty inherited list
eventually under the topological upper comparisons on the path; alternatively the strict stack inequalities
give the stability, since limiting equality of an original rim label with a loss boundary would be an actual
contact. Nonempty open intervals and switches of ordinary targets transmit stably by the strict Cauchy
record rule. Every further ordinary target prefix through strict preterminal observations lies also on paths
to an open interval of targets, and is read by common-prefix transmission from the dense lists. Thus the
resulting measured tree need not keep any priority of those lists. In particular auxiliary ordinary targets
added later by continuous fraction densities do not call for a new continuum disintegration.

The mesh for the one-branch-at-a-time countable encoding still shrinks. An infinite negative chain staying
in one caller misses at least the first \(N\) list targets after \(N\) switches (the first eligible is
strictly lost at each switch); all macroscopically passing leftover components in that caller stabilize by
finite targets as in Proposition~\ref{wall:returns}. Infinite chains containing repeated canonical restarts instead shrink by half-layer
active mesh or by nested \(O\) mesh of Lemma~\ref{wall:mesh} and compact disk containment. Siblings at any one path shrink by the
single-branch rule. This gives the full uniform shrinkage. Volume and spatial coupling for these routes now
use Proposition~\ref{cut:disk-spatial} and the finite-depth area dominance, since offspring interiors are area-additive and
the false-pinch and joining comparisons are along countably many compatible cut paths with this mesh. One
can accordingly extract universal observations using this countable primary encoding and then sample the
remaining embedded field/tree data from their actual continuum conditional law given \({\cal T}\),
independently of extra lattice readouts given that data input. This does not ask for determinism of the
fields or embedding from \({\cal T}\).

\end{proof}

\begin{definition}[Buffered observation arrays]\label{kern:buffered-observations}

Here are useful precise conventions for passage observations. One may keep nested countable triangle-cover
tests in abstract exploration coordinates before converting them to embedded tests. For example split
central and large positive rings into relative circumference chunks; split preterminal paths in compact
positive-length regimes into time blocks, designating separately large jumps and their fillings; keep also
terminal unswept fillings. A block includes all nondesignated jump contents born during it and the triangles
on rings produced there. Refine through finitely many paths and cells at a time, using continuity cutoffs
and then exhaustion. Each such piece stands for its Tutte triangles, hence for their flags with incidence
endpoints. Their upper images in any corresponding continuum realization are in the respective closed
swept/cut portions up to vanishing error. Meshes of such finite partitions can tend to zero: split all large
rings, isolate the large jumps, cut short ordinary segments and tiny terminal tails, and descend the
finitely passing cells. Thus with compact/open loss in space these cover \emph{all} lifts.

Keep numerical
path and triangle variational observations jointly for countable finite alphabets of such piece tests, in
any deterministic cost units under consideration (compactify nonnegative costs at infinity). For paths one
can keep attainability of ordered marked-point, range/confinement and upper cost requirements, including on
subpaths or finitely many paths at once, or equivalent hyperspace closed records; strict budget thresholds
can be approached from either side. Polyhedral extremal-length tests use families of curves with port and
confinement conditions specified by unions of flag pieces, enlarged to stars or inset as necessary. Exact
local oriented incidences and full flags with their fixed geometry suffice to evaluate such tests; star
completeness is needed only inside an intervening buffer.

Winding or transverse passages can equivalently
use strict ordered port and tube margins before and after testing Jordan annuli/quadrilaterals. All limiting
spatial readings below are with arbitrarily small port/confinement and budget slacks, or regularized over
intermediate open sets, not an identification of a microscopic EL at a limiting contour equality.

\end{definition}

\begin{lemma}[Equivalent refining arrays]\label{kern:equivalent-arrays}
In one common spatial realization, two refining piece arrays with vanishing upper-image mesh
and all-lift coverage determine the same buffered passage and variational profiles, using compact
constraints, open enlargements, and one-sided budget margins.
\end{lemma}
\begin{proof}

An equivalent countable navigation refined within matched raw pieces may be used to recover these buffered
spatial readings. Indeed any two finite systems with sufficiently small upper-image mesh in a \emph{common
coupled} realization impose the same inclusions from strict compact constraints to enlarged open constraints
by all-lift coverage under the spatial comparison. This works for paths, curves and all flags meeting a
prescribed protected part, and then for comparisons of budgets/variational costs by inclusions; it never
requires exact membership agreement at cover boundaries. Thus one can first extract intrinsic closed passage
profiles/witnesses with a primary alphabet, apply the conditional continuum realization as above, and
identify their spatial regularizations also with extra refining arrays jointly extracted when doing a local
calculation. Subsequence witnesses/attainability can equivalently diagonalize all positive margins. Only
such compatible localized observations, not conditional independence of arbitrary exterior-indexed
microscopic cover tests, are intended.

\end{proof}

\subsection{Finite mosaics and exact slot coordinates}

\begin{definition}[Finite mosaic]\label{kern:mosaic}

We now describe finite \emph{mosaics}. Add finitely many ordinary fresh target tests (and keep the sources)
at visited callers, routed with common prefixes. Stop before terminal, with all traversed side lengths in
positive compact ranges apart from designated jump transfers. Use:
\begin{itemize}
\item ordinary positive-duration sweeps with iid step kernel \eqref{eq:original-12};
\item individually designated macroscopic positive/negative jumps between sweeps, with both resulting lengths/specified child perimeters interior as needed; switch additional targets into detached disks in strict open lists, or schedule new tested callers after full canonical births;
\item unswept terminals including remaining polygons at stops. No second tested branch uses a terminal in its interior on this finite chart.
\end{itemize}
Divide each designated ring including the central ring by smooth cuts of its \textbf{outer} circumference
into positive macroscopic chunks. The entire construction is finite and includes signs, orientations, fork
and target instructions. Choices can use auxiliary marks with smooth positive densities: target marks, cut
positions, times, and designated jump observations by counting (time and size windows). All sweeps between
events have positive duration. Thus lattice multiplication/selection factors on a localized chart have
ordinary macro densities, including a factor of order \(m_P\) when counting each designated jump time (use
ordered counting and interior cutoffs). Equivalently integrate over adjacent sweep step counts as free
elapsed times when marking jumps. Exact survivor/eligibility restrictions are imposed; comparisons exhaust
by positive margins and by the stable inherited decisions above. Cuts at ordinary times can use
positive-density Poisson clocks or joint duration densities.

In the limit the added chart variables are
marked from \({\cal T}\) without extra passage conditioning: smooth independent marks/fractions and counts
with continuity restrictions transfer by common-prefix convergence, \(J_1\) large-jump matching and
ordinary-position convergence. When using counting, localize also the finite multiplicities, a tilt by input
data. Versions can thus be tested on positive-support charts and chosen with margins after the primary
extraction, not by making microscopic cost-based decisions.

\end{definition}

\begin{lemma}[Separated retention]\label{kern:retention}
For finitely many separated compact observations, finite mosaics of positive support can retain
all pieces and incidence stars needed by each observation inside its private buffer. Nuisance
positive-volume pieces away from the tests can be reserved separately.
\end{lemma}
\begin{proof}

For separated compact observations retain whole small terminal disks and whole ordinary sweeps (including
all their nondesignated jump contents) when needed there, and retain ring chunks and designated point
gluings locally. Pieces retained for one patch are entirely within its larger private buffer. Other ordinary
sweeps are away from the inner tests and all their stars with margin there. This can be arranged on a given
drawing: reach each of the finitely many canonical macroscopic cells still requiring traversal using finite
branching stabilization and recursion; choose sufficiently many further targets that uncharted pockets no
longer pass macroscopically across the buffers. Truncate the tested paths sufficiently late before each
target after all needed forks/discoveries, designate all crossing jumps on them, and take short time cuts
between, using continuity of trunks and vanishing diameters of other jumps/tails. These are positive-support
choices, with open target intervals to capture any demanded preterminal path prefix or fork. For two-target
forks into both sides require interior representatives there. Unseen macroscopic positive-volume pieces away
from tested data can be reserved separately for total-volume smoothing. No deterministic bound on a chart
size in this first simultaneous separated test is necessary.

\end{proof}

\paragraph{The exact pending-slot operation.}

There is an exact pending-slot description. Each cycle consists of boundary slots awaiting gluing. The
central ring produces its outer and inner cycles. At a designated face, consume the active slot against its
matching outer base on the new ring, emitting the other outer bases into the main frontier and the inner
cycle for its child. At a designated bridge, reverse-glue the two indicated slots, with all other slots
transported unchanged to the two components. Ring binary chunks contain the triangles/gluings along those
bases. Terminal disks consume their whole cyclic lists.

Over an ordinary sweep, the initial target and
ordinary prefixes through heights at or below each side's minimum propagate unchanged. Consume the active
and suffix slots strictly above those minima, and emit the final active and final suffixes above the same
untouched prefixes. Include here a possible artificial pass-through (e.g. final active reassigned from one
of those input suffix slots) recorded by the raw word itself.

Thus the whole local cobordism with its
ordered gluing tests, exposed input/output suffix origins relative to the initial ordinary tops, and all
nondesignated fillings can be generated without knowing deeper stack offsets. This is just the two-list step
rule with index heights \(1,\dots,L\) and \(1,\dots,R\) for ordinary slots. In nuisance sweeps pending slots
far from consumption are simply transported. Iteration of actual slot gluings and the raw pieces
reconstructs full triangle incidences; at an inner-protected vertex all flags in its star are local by the
topological upper comparison, so that cyclic/vertex identifications there need no unobserved extra exterior
vertex equivalences.

\begin{definition}[Independent-scale coordinates]\label{kern:coordinates}

We will use the following independent-scale lattice coordinates, in addition to elapsed sweep counts:
\begin{itemize}
\item \(n_a\), the number of outer bases of every chunk on every designated ring (their sum is the jump's outer \(k\)); and \(e_a=m_a-n_a\) where \(m_a\) is the inner count of that chunk in matching interleaving order;
\item the intervening ordinary-slot count \(j\) at each designated bridge in the indicated negative branch convention;
\item fresh mark indices from prescribed cyclic origins when new marks are made (sampled before their routes; inherited marks are transmitted instead);
\item both ordinary-length changes \(d\) over every sweep (keep sides also at steps of zero change).
\end{itemize}

Write \(Y\) for the first-scale list of all these integer coordinates \emph{except} \(e\), without yet
dividing by \(P\). Macroscopically \(e/P\to0\); one must \textbf{not} replace it by zero in the gluing
tests. At rings use a convention starting on a specified outer base with independent geometric \((1/2)\)
inner runs along successive outer bases, before smooth tilts. In the central mass this convention uses
\(\binom{k+l}{k}/(k+l)=\binom{k+l-1}{l}/k\), choosing a typical first outer base for the phase. Root/FK
rooting when present for counting does not affect inner tests. Normalizers and tilts by \(f_l\) at
designated faces, \eqref{eq:original-12} likelihood ratios, large-step masses times time and lattice units,
and mark/cut densities are smooth in the macroscopic parameters on compact interior ranges by the
established asymptotics; diffusive-scale \(e\) remains unconditioned by any exact circumference matching.

\end{definition}

\begin{lemma}[Endpoint local limit theorem and bridges]\label{kern:endpoint-llt}
Apart from the chunk discrepancies \(e\), the first-scale coordinates have the full integer lattice.
Here fluid endpoint ranges mean compact ranges of the rescaled endpoint coordinates,
which may have either sign, with durations bounded away from zero.
On compact positive-duration, fluid endpoint ranges, the ordinary increments have the joint uniform
Cauchy endpoint local limit theorem, and their \(J_1\) bridges converge conditionally on endpoints,
also when durations vary in a convergent positive compact range.
\end{lemma}
\begin{proof}

Other than the \(e\) coordinates, \(Y\) is on the full integer lattice with locally continuous bounding
densities/LLTs where needed. In particular positive \(d\)-duration gaps have the joint uniform Cauchy
endpoint LLT with \(J_1\) increment bridges converging conditionally on endpoints in compact fluid ranges.
Indeed \eqref{eq:original-12} has no displacement period (zero and unit generators), and convergence of
characteristic-function powers on compact frequency annuli at scale \(1/P\) bounds powers by an integrable
exponential of a smaller positive frequency power up to small fixed raw frequency, by dyadic annuli or the
tail equivalent \(w_{\pm k}\asymp(\log k)/k^2\). Away from zero powers decay exponentially. Fourier
inversion gives the assertion for endpoints, with strictly positive Cauchy densities. Restricting a bridge
to an initial fraction away from its end now uses bounded convergent likelihood ratios; reversal handles the
residual oscillation in a vanishing end fraction uniformly on endpoint compacts. This gives the conditional
bridge assertion, also for convergent positive durations on compacts. Attached marked Boltzmann fillings
continue to use the product rules along these paths. We will account separately for projection of this
lattice and for the conditional macroscopic tree data after chart smoothing.

\end{proof}

\subsection{Local projections and their integer lattices}

A raw retained piece records its internal geometry, increments, and running minima.
To assemble several pieces, we also need their relative slot origins. For a patch
$b$, the vector $O_b$ will contain retained sizes and sweep endpoint changes,
together with the integer translations needed for local gluing. For example, a
transported index $u$ may be tested against a raw running minimum $m$ by
\[
 u+a(Y,e)\le m.
\]
The affine integer expression $a(Y,e)$ belongs to the parameter vector; $m$
remains part of the retained raw word. The following proposition constructs the
required rows, including those whose limiting viable interval is a single point.

\begin{proposition}[Exact local gluing projections]\label{kern:local-projections}
Fix a finite mosaic history on a supported valid chart. Along a convergent lattice realization,
all sufficiently large approximations have their buffered local readouts determined by the retained
raw pieces and finitely many affine integer gluing tests. The tests retain every viable
closed macroscopic prefix, including singleton cases. Every binding translation is accounted for
by arcs inside the private buffer.
\end{proposition}
\begin{proof}\leavevmode

\paragraph{Interfaces and their origins.}

Here are details on choosing translation rows for such a mosaic. All charts can have extra spatial collars:
in discussing one inner test keep an enlarged compact protection and still larger open private buffers.
Sweeps or uncharted disks coming near enough to the protection are retained entire for that patch, with
closure in the private buffer; likewise ring chunks there. Include all portions needed on the two sides and
all star neighborhoods for the inner test. A sweep's exposed suffix arcs, including limit endpoints, are in
its closed swept/attached region: nonpersistent input labels are consumed toward the tip or into its
attached cuttings, and new labels come from its discoveries or the moving tips. This also follows by tracing
a.e. heights by strict stack transmission, then continuity. The regions here for ordinary sweeps need not
include anything in the interiors of \emph{designated} jump disks, which have separate pieces.

The following procedure can use redundant tests. List output interfaces (ring chunk side lists excluding
slots separately paired, sweep output active and suffixes) and input interfaces (sweep input active and
suffixes, point consumers at designated jumps with the indicated attachments, and terminal cycles). An
active used in two interfaces in succession counts by the corresponding pass-through conventions. Use
orientation of boundary of the \emph{unfilled} disk for indexing unchanged slots; thus interval transfers
preserve cyclic orientation. Origins for suffix labels can be taken at the corresponding active (use the
\textbf{final} top for an output sweep suffix). Active and other singleton slots can also be kept
separately. Ring chunk origins are at their own ends on each side. State root indices when routing can be
taken at actives or attachment points, using a rotation when making a fresh mark/activation. Fixed offsets
like using the first inner base after a starting outer cut use interleaving order (run counts from
successive outer bases, including runs of length zero); one can root the whole inner cycle at its first base
in this convention. Put all extra subdivisions at cuts in outer-base order, first-base conventions fixed
before the runs. For macro chunks here absence of inner bases is negligible.

Ring adjacency \emph{around}
the ring between chunks, with radial links, is deterministic along their adjacent ends; record those
instructions too. It is not tested by an equality of an outer and inner slot height.

\paragraph{Transport through unchanged slots.}

From a needed output slot, trace to its next consumption in the finite mosaic. This passes along unchanged
slots only. A chart of forked branches is ordered as successive polygon operations, doing each common prefix
once up to a designated split; no later operation before consumption looks inside another pending interval's
slot. In an ordinary sweep, the unchanged interval through the target (both ordinary prefixes and target
together) has a single translation, given by the side length changes. Its end tests use the two running
minima relative to the initial ordinary tops. It need \textbf{not} be subdivided at the target. In
particular many remote sweeps with the same possibly local target need not make wrap or point tests there.
Designated splits transfer contiguous intervals to the two polygons; positive births transfer the old
interval except the consumed slot to the post state and insert the designated outer list, keeping the inner
list for the new disk. Rotations at state changes, initiative changes (activation or redirection), wraps,
inclusion/exclusion of distinguished slots and final index comparisons just give orientation and affine
integer translations and inequalities. This is interval counting on cyclic lists with a slot as one unit,
not modulo a spatial boundary-identification relation. A birth-list label may in particular be tested
directly against point consumption there.

Every translation, and every point or interface threshold relative
to a state origin, is affine integer in (\(Y,e\)); the only additional comparisons at ordinary sweeps can
use their raw running minima with the translated index anchored at top. Indeed current cycle sizes and
marked positions propagate at a negative split by cutting out \(j\) ordinary slots and deleting the two
paired slots, with remaining positions rotated from the prescribed ends; at a designated face by adding
\(k-2\) slots to the parent and using \(\sum m_a\) as inner perimeter; over a sweep by adding its two net
ordinary changes, with the position of the target on the persistent interval unchanged relative to its
slots. Mark indices can first be created at birth of the cycle on which they are freshly chosen; comparisons
along a continuation instead transmit the old indices. No duration or branch extinction-time equality is
imposed.

\paragraph{Closed viable prefixes.}

In forming the tests consider all viable \textbf{closed macroscopic prefixes} of these finitely many
translation cases, testing heights in perimeter units and \(e/P=0\), allowing equality at endpoints, with
specified side/lift at wraps. Clip the source interfaces as needed by finitely many closed coordinate
intervals (rational bands with slack are possible) whose points on the valid interface are in the
protection, covering all required inner slots with margin. For a fixed history the viable source-coordinate
set is a closed interval or a point, by linear interval tests (raw minima now read at macroscopic values).
At each viable prefix the represented point is still the source's physical point: transmissions extend to
endpoints by the exact stack and aperture rule. Do the following:
\begin{itemize}
\item Drop cases with no viable prefix. Comparisons at the next step that cannot be equalities on this interval give just a pass/fail bit there.
\item At comparisons that can bind, record exact affine shifts from the source coordinate to the threshold. At a running-minimum test keep the shift to its deterministic top anchor, with the minimum then a raw comparison. The active/paired point conventions and strictness on the lattice are kept, not a free link at equality.
\item For passage to a final interface consumer record the full relative translation for index pairing there, including orientation. Use the local point-gluing instruction when applicable.
\end{itemize}

\paragraph{Sufficiency and incidence completeness.}

These rows suffice asymptotically exactly on the chart: by compactness and convergence, omitted strict
comparisons remain valid even for approaching compatible prefixes (or exclude the entire case); induct on
steps. Binding raw minimum comparisons all occur in retained sweeps, since equality places the threshold
point at a tip/limit there. Keep clipping margins so that all lifts near inner tests are covered. Conversely
a proposed gluing in the clipped tests cannot skip an omitted obstruction, even for a singleton limiting
case.

Rows for matching to a consumer or for a binding threshold have a useful macroscopic accounting:
evaluate at \emph{one} common viable point (the binding point when needed). The offset is a signed length
comparison along the source interface and the consumer arc or the arc to the threshold's deterministic
anchor there. Both live within the private buffer (at a minimum use the corresponding exposed arc up to the
minimum). At a designated point threshold one can use the threshold itself, not the length back to a distant
cycle origin. Thus absolute macroscopic translations for these rows use local lengths. Viability bits and
ordering are separate instructions. Tests that can occur in the inner part but originate on a chunk's or
operation's immediate neighbors can always be included with the same buffer.

With raw pieces and these integer parameters we can therefore compute passage readouts inside, using open
safety margins. All base edge attachments come from the consumed-slot gluings, including the reverse pairing
at bridges. Inside raw rings/cuttings and along retained neighboring ring chunks we have the full triangle
incidences. Vertex and star identifications are generated by following incident triangles through such
gluings: equivalence about a vertex in the final closed triangulation can be followed in its cyclic star,
\emph{all} projecting locally there by the incidence consistency. Singleton perimeter-zero and
boundary-pinched attachments inside a raw piece keep their true local equivalences. Thus, for instance,
paths using original primal edges (the indicated blue diagonals) and local polyhedral curves on the actual
flags can be used; this does not substitute graph paths on a boundary quotient made only at continuum
resolution.

\end{proof}

\begin{definition}[Projected local parameters]\label{kern:projected-parameters}

For patch \(b\) include in \(O_b\) the retained outer chunk sizes \(n_r\), retained sweep endpoints \(d_r\),
the foregoing required parameter rows, and entire retained terminal sizes. (A terminal includes any unswept
child filling after a designated operation.) An untested chunk strictly local can be included too. Keep
ordinary durations separately. We have for their concatenation
\begin{equation}\label{eq:original-14}
O=M Y + C e + c.
\end{equation}
Everything in this identity is unscaled. Use only fixed finite instruction histories when writing a matrix,
even if the choice among them is made by macroscopic comparisons.

\end{definition}

\begin{lemma}[Integral product splitting]\label{kern:integer-splitting}
For \(L=M\mathbb Z^{\dim Y}\), the image lattice in \eqref{eq:original-14} splits as
\(L=\prod_b L_b\). Remote discrepancy columns can be absorbed simultaneously by integer
translations in these factors. At fixed retained \(n_r,d_r\), the remaining fibers are product
affine lattices. Local discrepancy columns have representatives bounded by local complexity.
\end{lemma}
\begin{proof}\leavevmode

\paragraph{Paired atoms on primitive circles.}

We use an elementary algebraic criterion. If a generating family of
$\mathbb Z^{\dim Y}$ has images under $M$ each supported in just one patch,
then $M\mathbb Z^{\dim Y}=\prod_b M_b\mathbb Z^{\dim Y}$: each patch
projection can be formed independently by integer sums of these images.
We construct the generators by boundary-label counting. Think of adding a unit
at one generic point of a \textbf{primitive} circle
(central circle or positive aperture) on \emph{both} sides. Propagate its unit count in the same physical
intervals of slots at the macro states, keeping all marks, roots, cuts and pinch positions fixed. More
explicitly set \(\delta Y\) by counting this atom in designated chunks on the outer lists, in detached
pre-intervals of negative jumps, and in indexed arcs for mark positions, and by differences of its counts on
the two ordinary lists at sweep ends versus starts. It has \(\delta e=0\). This variation is defined in the
continuum interval arithmetic; it does not ask for insertion surgery on a lattice path with endpoints
constrained.

At the finitely many states, almost every boundary label is strictly inherited on an open interval from
an original boundary or a positive jump: trace backwards using the closed record range of the reversed
Cauchy paths, null in height. This applies at the ordinary continuous-density times and pre/post designated
jumps here by the Lévy jump counting law and absolute continuity on compact surviving pieces; through forks
and caller changes transmit previous intervals as well. A label on one such circle side occurs at most once
(apart from cut endpoints) in a cycle; its other side does not later merge into that cycle. Its appearances
in finitely many previously measured state intervals are locally constant off endpoints/null cut locations.
Indeed along a sweep the persistent interval truncates lists, the rest is new by backward tracing, and a
fixed positive aperture created there likewise transmits by prefix truncation. Across the finitely many
other operations these are just interval transfers. Thus generic choices can avoid specified ends, including
points selected to account for finitely many binding rows.

\paragraph{Locality of projected atom counts.}

The resulting \(M_b\delta Y\) counts the atom along just the local arc comparisons above (and local
size/sweep rows). Indeed all translations along unchanged intervals identify counts there, and one can sum
on either side of the common viable labels, even when using several successive closures instead of an open
overlap: intermediate coordinates are then the \emph{same interval labels} including indicated wrap lifts,
not arbitrary coincident physical points. Endpoint positions carry none of the generic atom. At a local
sweep \(d_r\) uses differences with the whole unchanged prefix cancelling. This proves the assertion by
addition/subtraction of the indexed arcs, or iteration of the cycle updates defining \eqref{eq:original-14}.
Each entry has bounded magnitude by a fixed number of incidence counts at the two local observations,
independent of the length of the transfer history. In particular \(M\delta Y\) is supported in at most one
private buffer.

\paragraph{Triangular unit generators and discrepancy columns.}

The group generated by the paired \(\delta Y\)'s is the \textbf{whole} parameter lattice for \(Y\):
\begin{itemize}
\item At birth of each designated outer circle take a generic location in each outer chunk. These are unit variations of their sizes among the coordinates created there and earlier.
\item At a negative jump straddle its lower loss threshold by two such placements on the same inherited open interval, subtracting the placements. This varies \(j\) by one, cancelling all earlier variables. Indeed the jump threshold has density in the pre-list at a designated loss (use Lévy jump counting before any restrictions). Fresh source/target marks, indexed by length, allow the same straddling for their individual positions.
\item At the end (or just before the designated end jump) of each positive-duration sweep, on \textbf{each} side the end value strictly exceeds the minimum almost surely, by reversed Cauchy oscillation at ordinary or jump-counted times. Thus a positive fraction of heights there traces back to new nondesignated positive insertions during this sweep. Take an atom on such an exposed interval; its paired-side copy is in the jump child. This varies the corresponding \(d\) by one without varying the other endpoint change or earlier variables.
\end{itemize}

Order coordinates by the finite branch scheduling, measuring marks and sizes where first introduced, not at
a later fork. These are triangular unit generators. The genericity assertions hold on actual supported
charts (and unit-area integrated-perimeter versions by the scale Palm comparison); in particular they are
not claims at adversarial exact Cauchy terminal times. Consequently the \textbf{integer} image \(L=M\mathbb
Z^{\dim Y}\) splits as the product of its patch images \(L_b\), since it is generated by projected atoms
singly supported.

For a column \(a\) of \(C\) now use a placement on the corresponding designated chunk's
inner side only (\(\delta e_a=1\)). The same counting holds for \(M\delta Y+C_a\); translation routes do not
equate outer and inner indices through an interleaving, which is kept as a raw chunk. Cuts stay at the
paired physical endpoints, with counts allowed to vary unequally. Thus \(C_{b,a}\in L_b\) whenever that
chunk has a generic location off the patch's length supports. Make all other such chunks retained for that
patch (they are inside its private buffer). For local columns the calculation still gives bounded
representatives modulo \(L_b\). Remote translations of columns can now be absorbed \textbf{simultaneously}
integrally by product splitting. Projection also onto the \(n_r,d_r\) rows causes no new condition—they are
included in this proof and those first coordinates themselves are free. Fibers for the other local rows at
fixed first coordinates are product affine lattices.\qedhere

\end{proof}

\begin{lemma}[Bounded local lattice and residue types]\label{kern:finite-lattice-types}
Fix a bound on the number of local retained operations, arc comparisons, and distinguished-slot
conventions, with their supports in a strict inner buffer. The possible image lattices \(L_b\),
local discrepancy-column representatives, and constant classes \(c_b\bmod L_b\) form a finite
collection, even when remote transport histories have unbounded length.
\end{lemma}
\begin{proof}\leavevmode

\paragraph{Weighted actual slots.}

To control the constant term, we construct an integer vector $Y(w)$ for which
$M_bY(w)+c_b$ is bounded solely by local complexity. Since $M_bY(w)\in L_b$,
this bounds the possible residue classes of $c_b$. We use weighted counts of
actual slots. Suppose the entries of \(O_b\),
contributing local operations/marked slots and subdivisions are bounded in number, with all arc supports for
these entries in a strict inner part of the private neighborhood; there need be no bound on remote
transports. Use an actual lattice realization along a convergent chart approximation (available at the
typical input points, also by the trajectory couplings). Give integer weights to actual slots created on
circle sides and carried unchanged before consumption. Put 1 on all distinguished slots used at finite state
ends, marks, origins and cut or point-operation conventions (include their boundedly many immediate endpoint
slots when used in unit offsets), zero initially on the rest. It is not necessary to do this at internal
steps of an ordinary sweep. Adjust by signed weights at unconstrained slots within each designated chunk so
its two weighted totals agree.

Near all arc comparisons for \(b\) the total absolute slot weight used can be
bounded just by local complexity. Indeed distant distinguished slots do not appear there; if a chunk has
distant weighted slots and an unbounded possible mismatch, its adjustment can be made on an open subinterval
outside the local supports, avoiding the finitely many constrained slots of that realization. Use buffer
margins in this argument (bound distinguished slots on a fixed enlargement of the arc supports). Repeated
naming of an unchanged slot, e.g. a target, takes just one weight. Here weights are on actual side-slots,
not simultaneously on their glued opposite slots, except for imposed unit conventions.

\paragraph{The unit constants and closed-endpoint accounting.}

With weighted sizes, changes and marks \(Y(w)\) computed as before on the actual cycles, formulas for the
shifts have the \textbf{same unit terms}, and \(e(w)=0\). They result by adding whole interval counts except
for the indicated distinguished endpoint conventions. Thus the values in the weighted arithmetic are
\(M_bY(w)+c_b\). They are bounded by local weighted counts and a bounded number of local unit conventions.

To check this even for merely closed viability, index by integer gaps before successive slots, in the
orientation preserved along an unchanged interval, extending cyclic lifts by their periods. For each
primitive unchanged-interval arrow through a common viable label one can take strictly transmitted gaps
approaching that label; the translation formula, with the indicated full turns added if necessary, sends
them to lifted gaps approaching the specified end coordinate. Indeed at a bridge these are the preserved
sublists in one of the two positive macroscopic cycles; at a positive birth the inserted and preserved arcs
have prescribed starts/ends; over a sweep use the persistent interval through the target, of positive size
with approach allowed up to its endpoints. Pure wraps or reindexings just take cyclic lifts. A singleton
point instruction can instead use its very slot or endpoint. On actual transferred gaps the weighted shift
equation is exact in these lifts, even if one must evaluate a few \(o(P)\) positions beyond a stated wrap
range; all involved gaps near the endpoint represent positions near the protected point by contour
continuity.

Approximating gaps for consecutive arrows may disagree by \(o(P)\) slots at that same lifted
coordinate. The weighted difference then costs only local absolute weights. Moreover changes of
approximating gap are needed only at local threshold operations: throughout any strict passage along remote
arrows a single approaching slot actually transfers unchanged with the indicated lifts. Thus the error in
summing the shift comparisons this way is bounded by local complexity, also if no single lattice slot makes
the whole proposed passage. At a raw-minimum test one only counts back on its observing arc to the top
anchor; the minimum itself need not be a distinguished weighted slot. Orientation conversions at local
readout interfaces can similarly use their endpoint unit conventions or the bounded local weighted errors.
This proves bounded representability of \(c_b\) modulo \(L_b\).

\paragraph{Eliminating the remote history.}

In this argument remote transports can be taken rooted at remote actives until they actually switch to a
local origin. The persistent interval through a target during many such sweeps makes no subdivision at that
target. Thus, when bounding local complexity, no error or threshold cost proportional to that remote count
is incurred. More generally a remote operation none of whose interval ends or point choices bind on a given
clipped viable range transports that entire connected range uniformly. Equality would put one such end at
its protected points. Hence repeated subdivisions and choices of offsets are bounded in number using only
the operations with ends there, including those for matching ring chunks and activation, and finite branch
tags for the resulting chains. For fixed local complexity the possibilities for \(L_b\), the remaining
column representatives and the residual class of \(c_b\) are now finite: \(L_b\) itself is generated by
bounded local atom-count vectors in bounded dimension. This finiteness does not yet assert that there is
bounded local chart complexity over conditional extensions; we address the chart choice below.

\end{proof}

\subsection{Reference experiments and the whole-tree product rule}

We now pass from the integer description to conditional laws. Each retained
block $b$ has four kinds of data:
\begin{itemize}
\item $S_b$ consists of the limiting projected parameters $O_b/P$, retained
  durations, and macroscopic increment paths of the retained sweeps;
\item $U_b^*$ contains the limiting raw-piece observations and their countable
  continuation tests, all extracted on the same underlying raw maps;
\item $I_b$ specifies the entry modes and transmitted boundary intervals supplied
  by the rest of the chart;
\item $G_b=g_b(U_b^*,I_b)$ is the actual continuation tree, with volume labels,
  obtained by following those entry instructions.
\end{itemize}
The augmentation and the measurable projection $g_b$ are constructed below.
In particular, $G_b$ uses the actual prescribed or fresh source at each entry;
it does not reveal every hypothetical continuation retained in $U_b^*$.
The full chart data $S$ include all the $S_b$ and the nuisance parameters.
The following elementary fact explains the probability argument.

\begin{lemma}[Conditioning independent blocks]\label{kern:conditional-blocks}
Let all spaces be standard Borel. Start with independent pairs $(S_b,U_b)$
of laws $\mu_b$, $1\le b\le B$, and sample a completion $S_0$ from a kernel
$q(\mathrm ds_0\mid S_1,\ldots,S_B)$. Put $S=(S_0,S_1,\ldots,S_B)$.
One may reweight this law by a nonnegative function of $S$ with finite positive
integral and add nuisance observations whose conditional law given $(S,U_1,\ldots,U_B)$
depends only on $S$. For measurable $I_b=I_b(S)$ and
$G_b=g_b(U_b,I_b)$, the conditional law of the $U_b$, given $S$, the nuisance
observations, and all the $G_b$, is the product of their separate reference
conditional laws given $(S_b,I_b,G_b)$.
\end{lemma}
\begin{proof}
Before the separate projections are revealed, the conditional law given $S$
is $\prod_b\mu_b(\mathrm du_b\mid S_b)$. Reweighting by a function of $S$
and adding the stated nuisance observations do not change it. For fixed $S$,
each $G_b$ is a function of $U_b$ alone, with fixed parameter $I_b$.
Disintegrating over these separate functions gives the assertion. Measurable
versions in $I_b$ can be obtained by countable refining partitions of the
projection range, as in the application below.
\end{proof}

We must implement these hypotheses for the finite maps: construct independent
references, prove the stated dependence of the density, and show that the
actual projection trees reconstruct the full primary input.


\begin{definition}[Independent block reference experiments]\label{kern:reference}

Here are the reference experiments in more detail. They can use, for example, smooth everywhere positive
lattice densities by scaled Gaussians (one-sided densities for positive free coordinates); assembly and weak
density comparisons need only be made on compact fluid ranges. Dummy outputs can be used for invalid
parameters. All cost units and exact lattice readouts in a kernel comparison use the same sequence;
concurrent extractions below, including references, always allow passing to further subsequences.

Take independently for each local instruction block \(b\) (the union of pieces kept for that patch):
\begin{enumerate}
\item Retained ordinary durations and retained \(n_r\) by positive lattice densities at \(m_P,P\). Retained raw words by the iid \eqref{eq:original-12} scheme with all nondesignated ring/filling contents. Thus positive ordinary steps sample interleavings and fillings with their true step tilts; negative steps include the detached filling, not the new target polygon. Keep the ordinary lists relative to tops without a bottom obstruction in this raw generation. Independently, for the retained chunks of designated rings use geometric runs given their outer sizes, producing \(e_r\) exactly, before any total inner filling tilt.
\item The other \(O_b\) parameters on the affine lattice fiber from \eqref{eq:original-14} over \(n_r,d_r\), translated by the local columns on \(e_r\) and the local constant class, with positive smooth lattice density at scale \(P\). Columns/classes here are modulo \(L_b\); remote columns are dropped. Generate retained terminal fillings with their true disk laws at the resulting sizes, independently given parameters.
\end{enumerate}

The choices include side colors, orientations and discrete point conventions as instructions; costs use the
actual fixed flags/primal incidence, not a rotation-averaged metric. Variation of smooth densities or fiber
sections here is harmless if the ratios limit to functions of conditioned macroscopic data. Fibers have
constant lattice covolumes; whenever their constrained inputs after scaling converge and \(e/P\to 0\),
normalized Gaussian choices converge by ordinary lattice sums (affine translations can be recentered within
bounded lattice error on the moving plane). Include all desired partial-gluing, triangle/curve and primal
path observations using the exact integer readouts, \textbf{before} taking limits.

\end{definition}

\begin{lemma}[Actual continuation projections]\label{kern:whole-tree}
The retained block can be augmented to a countably encoded payload \(U_b^*\), with initial
conditioning parameters \(S_b\). Given its entry modes and transmitted intervals \(I_b\), its
actual continuation tree is a measurable projection \(G_b=g_b(U_b^*,I_b)\). The chart axes and
these projections for all local and nuisance entries reconstruct the whole measured tree
\(\mathcal T\), including volume labels.
\end{lemma}
\begin{proof}\leavevmode

\paragraph{Entries and inherited intervals.}

Call an uncharted entry a macroscopic attached positive or negative disk within an ordinary word, or an
unswept terminal \(T_s\). Its cycle can be rooted at a prescribed active or attachment slot. At negative
cuttings or in a post polygon still serving the original target experiment, let \(I\) be the transmitted
interval of original caller target labels on this cycle, in its relative-length coordinates. If nonempty the
future parent paths there are exactly continuations from the active/pinch aimed at that open interval (up to
null/endpoint labels), not explorations from a newly resampled origin. If empty, or for a separate positive
disk, there is a fresh start. Thereafter apply the same rule recursively. In integer arithmetic the nonempty
ordinary positions in \(I\) transmit by contiguous slot counts; endpoints at special actives need only be
considered through their strict-margin limits.

Initial \(I\)'s converge and are known macroscopically from
the chart lengths/paths: along a path to a still transmitted original target the available ordinary labels
on each side from that target are kept by prefix truncation, through separation by the bridge rules.
Nonempty open portions persist. Empty-list restart decisions at positive-size negatives are stable: equality
with an original untouched threshold would put the pinch on the old boundary by horizontal transmission; an
empty-interval complete negative is a compactly contained return. Equivalently for finite continuations,
thresholds of an initial eligible interval and previous surviving minima are a.s. not hit exactly by a
negative jump landing, by localization and Lévy jump counting. Compact survival up to stops to an original
target keeps a nonempty interval at the stop. At deeper tests these same arguments apply through converging
strict transmitted labels.

\paragraph{Augmentation on the same raw map.}

Augment local payloads as follows, so as to be able to condition on the \textbf{whole macroscopic tree}.
Within each such raw filling compare both the possible prescribed-active and a possible fresh-source start
if needed, testing paths to dense ordinary boundary fractions in its rooted coordinates. Within jump
offspring make the corresponding tests recursively (transmitted continuation active or independent fresh
assignment). These are on the raw map with its signs, with independent uniform fresh-source slots; use a
common assignment at a polygon birth shared by compatible prefixes.

More formally on the lattice one can
assign independent slots to possible cut occurrences keyed by operation histories, sharing assignments
before targets actually diverge. The steps before a split are independent of which still admitted target is
followed (at the splitting bridge each list's continuation has the prescribed new active on it). If a birth
uses a fresh source in a compatible actual all-slot computation, all routes rediscovering that birth share
its operative history. Hence these extra assignments can include the actual ones without bias; whether one
uses an assignment as a new independent source is decided without looking at it. Paths with conflicting
hypothetical source choices may be tested jointly on the \emph{same} raw map; no product disk/branch limit
for the conflicting pair is being imposed.

\paragraph{Extraction and initial conditioning.}

Use countably encoded compact-path, jump and cell-volume labels of these augmentations and refining piece
tests/readouts along them, as well as the retained step paths. Extract limits with the local parameters,
independently across \(b\) in product experiments. Jump entries and recursive tests can be indexed by
ordinary continuity cutoffs, length order and surviving pieces; any one consistent ancestral test in a fresh
filling has the compact convergences already proved. For unused tests one may just use subsequential arrays
(or compactified outputs), not identify their joint law.

Thus disintegrations can be fixed for the countable
lattice types, finite readout alphabets and refinements. Write \(U_b^*\) for such local limiting payload
data, and \(S_b\) for its initial conditioned parameters (\(O_b/P\) in the limit, retained durations in time
units, macroscopic increment paths of retained chart sweeps). The local kernels will first condition on
\(S_b\). With entry information \(I_b\) (intervals/root modes for this block's uncharted entries) they
condition further on only the \textbf{actual projected} continuation trees \(G_b=g_b(U_b^*,I_b)\), with
volumes, not on every augmentation.

\paragraph{Projection and reconstruction of the whole input.}

To spell out that projection: at an entry with a transmitted open interval use the dense tests strictly
inside it. Include positive offspring with fresh starts. At a negative offspring continue into the
transmitted open interval there if present, again using labels strictly inside (one can switch to dense
labels in that entry's own coordinates); if none, use the fresh assignment and a new full caller. Prescribed
directions/coordinate orientations are part of the entry conventions. In a convergence coupling of true
compatible data, every fixed eligible augmented compact test is actually present in the all-slot tree
eventually, with the same pre-separation paths and switch into a loss interior. A target chosen in such a
child inside the transmission range would have shared the prefix to its birth.

Conversely all ordinary
compact preterminal path observations are read from these lists: take a target from the dense tests inside
the open interval still sharing the desired prefix (strictly so with room), then use the known coordinate
changes. The same holds on the lattice eventually up to the compared preterminal observations, by strict
stack comparisons. Apply successively at finite matched depths; canonical source choices there are common
when used. Joint convergence of these finite compatible observations, including volumes, is just the
route convergence of Propositions~\ref{cut:marked-targets} and~\ref{cut:joint-volumes}: the new tests in an uncharted entry, selected by its incoming lengths/labels
before looking inside it, can use arbitrary converging ordinary positions strictly within the admitted arcs.
At deeper selections exact emptiness and compact tests are used as above. This proves the assertions a.s. on
actual supports by countable tests, avoiding any continuity demand on selecting a target exactly at an
interval end. Thus \(g_b\) can use open-prefix projection on the extracted arrays (Borel fallback
definitions elsewhere). Given chart data this projected tree is read from \({\cal T}\) as well.

Adding chart
axis data and these projection trees for \textbf{all} uncharted entries, local and nuisance, conversely
specifies \({\cal T}\): another branch either still shares the path, splits into an attached negative,
continues in a stopped cycle, or is a restart; label rotations and old caller identities are known from the
transmissions. Volume labels for parent parts are accounted for by area additivity (axis skeletons
themselves have zero residual area). All assertions need only the measured unranked positive-margin tree.
Extra primary target lists if randomized can be drawn separately; they do not encode alternative sources
when an actual inherited interval is used.

\end{proof}

\begin{lemma}[Chart density comparison]\label{kern:chart-density}
On compact valid chart tubes, the actual marked chart law is a bounded convergent change of
density from the product reference experiment completed by nuisance data. Its limiting density
depends only on the full chart scales and increment paths, and on nuisance discrepancy variables;
it does not depend on the retained passage readouts. These tubes exhaust the actual marked law.
\end{lemma}
\begin{proof}\leavevmode

\paragraph{Completing the exact integer fibers.}

We verify the chart density comparison giving this prescription. First take unconditioned pair bases with a
smooth positive scale density on the first perimeter in \(P\) units, or the unconditioned bases \(b_n\) on
interior perimeter windows. Given the local reference outputs, freely sample other durations, a lattice
Gaussian at scale \(\sqrt P\) for all nuisance chunk discrepancies \(e_o\), then the remaining exact
parameters \(Y\) on the completion fiber of \eqref{eq:original-14} with positive Gaussian density at scale
\(P\). Existence is exact by product splitting including the \(n_r,d_r\) rows. Translations for the
\(e_o\)'s and cross-columns are integer translations, \(o(P)\) on typical compacts here.

On integration, the
limiting conditional completion parameters other than \(e_o/\sqrt P\) have laws depending only on the
first-scale conditioned parameters, and the latter discrepancy draw is independent Gaussian. For the
nuisance chunks fill binary orders conditionally on totals. For nuisance ordinary paths use iid increments
bridged to the sampled endpoints, with fillings as at ordinary steps, and for nuisance terminals use fresh
fillings. On fluid ranges the bridge limit is continuous conditionally on the initial parameters by the
\eqref{eq:original-12} LLT argument. Nuisance macroscopic increment laws added this way thus need no
conditioning on other local payloads. Binary chunk orders align the two sides asymptotically as before. The
nuisance discrepancy Gaussians need only correct the chunk sum likelihoods.

\paragraph{The density factors.}

Indeed the density for the \emph{actual marked chart} relative to this generation has macroscopic limits on
interior valid tubes:
\begin{itemize}
\item Each designated negative or positive outer jump size costs a lattice mass of order \(1/(m_P P)\) with smooth scaled equivalent, and its free counting time contributes \(m_P\); other time cuts can sample with smooth duration densities. Several outer chunk sizes in a positive jump or central ring use a total-size density and smooth subdivisions, yielding full-dimensional density. Uniform marks sample their first indices with continuous scaled densities as well.
\item Active-path likelihoods \(f_{\rm end}/f_{\rm start}\), where the indices are full perimeter counts, and other disk/ring factors are the \eqref{eq:original-5}, \eqref{eq:original-12} rules. At designated faces the tilt of geometric runs by \(f_l\) has continuous limiting ratio on compact fluid ranges with \(l/k\to 1\); at the central ring the outer-rooted convention above and the perimeter equivalents do the same. There is no exact common-circumference bridge. For nuisance chunks the pre-tilt negative-binomial sum mass has the Gaussian lattice equivalent on \(|e_a|=O(\sqrt P)\), \(n_a/P\) in positive compacts, giving a limiting ratio depending on the Gaussian coordinate and \(n_a/P\). For retained chunks use the exact run sampling in the reference instead.
\item Counting \(Y\) by first projecting \(O\), then the completion fiber, just uses fixed integer-lattice covolumes/scaling powers. The \(n_r,d_r\) coordinates are free first coordinates; for \(d_r\) use the exact unconditional raw sampling from the reference. For preassigned nuisance sweep endpoints \(d_o\) the iid bridge LLT supplies the continuous joint endpoint densities. Thus all powers cancel as for full lattice densities of the other \(Y\)'s at \(P\) scale. Explicitly at fixed \(n_r,d_r,e\), projection and completion sum to one summation over the other \(Y\)'s, whose dimension is the sum of the projected free ranks and the completion rank. Positive smooth fiber mass choices cost the inverse of those lattice powers with smooth limiting factors. Marked jump counts use ordered times among distinct operations, not an equality between two free times of retraced common prefixes.
\end{itemize}

\paragraph{Localization and exhaustion.}

These products follow by sequential peeling with fresh conditional polygons. Admissibility and fork
instructions only restrict them. One may localize by continuity cutoffs to bounded ranges with side lengths
positive throughout sweeps, jump sizes, used terminal sizes and durations interior, marked-target
comparisons and required transmissions strict, and discrepancy tails bounded (or truncate the nuisance
Gaussian ratio then exhaust). Survival and fork/restart requirements for the finite chart exhaust by such
stable tests; where a duration ends before a fork, still transmitted ordinary targets have positive
clearance, and at designated separations use strict jump interiors. Common slot/end conventions themselves
are imposed exactly, not as additional density restrictions.

The perimeter equivalents and Riemann-sum and
lattice bounds above give bounded convergent ratios on these restrictions, depending just on the stated
scales/increment data (and the nuisance discrepancy coordinates), not on passage bits. Cutoffs can avoid
weak boundaries; the first-scale parameter law before admissibility has the usual positive density class
with iid Cauchy paths (the completion fiber choice restores full density over the unconstrained coordinates
conditional on retained endpoints, and nuisance bridges have their positive endpoint densities).
Alternatively use continuous localization supported strictly on valid tubes. Exhaustion in the actual marked
law uses the path couplings and asymptotic tightness of discrepancies by geometric sums and smooth tilts.
Deterministic overall factors, such as finite-mass rather than probability chart conventions, do not affect
the conditioned claim.

Here is the resulting weak-limit identity. Fix one continuity localization.
Write $\mathsf R_P$ for the unlocalized completed reference measure and
$\mathsf A_P$ for the actual marked-chart measure restricted to this
localization. All chart restrictions are included in the density $w_P$;
the reference retains its independent-block and completion construction. The exact peeling calculation gives
\[
 \mathsf A_P(dv)=w_P(v)\mathsf R_P(dv).
\]
Let $U_P$ denote the countably encoded augmented observations and readouts
extracted from the raw blocks, and put $Z_P=e_o/\sqrt P$ for the nuisance
discrepancies. Along their common joint extraction, the bounded ratio
calculation above gives
\[
 (S_P,Z_P,U_P,w_P)\ \xrightarrow{\ d\ }\ (S,Z,U,w(S,Z)).
\]
Indeed each density factor converges in the same path/parameter coupling,
and the localization indicators converge at their continuity boundaries.
The $w_P$ are uniformly bounded on the chosen localization. Thus, for every
bounded continuous probe $F$ of the retained arrays and parameters,
\[
 \lim_{P\to\infty}\int F(S_P,Z_P,U_P)\,d\mathsf A_P
   =\int F(S,Z,U)w(S,Z)\,d\mathsf R.
\]
The limiting reference generation samples its completion and $Z$ independently
of the block payloads given the block parameters. Integrating $Z$ therefore
leaves a density depending only on $S$. This identity, followed by exhaustion,
is the change-of-measure statement used below.

\end{proof}

\begin{theorem}[Product kernels given the whole tree]\label{kern:product-kernels}
Consider any compatible joint extraction of countably many local readouts from separated retained
blocks. There are reference disintegrations \(k_b\), fixed for the countable local types, such that
the conditional law of the retained payloads given \((\mathcal T,\text{chart})\) is
\begin{equation}\label{eq:original-15}
 \prod_b k_b(dU_b^*\mid S_b,I_b,G_b).
\end{equation}
The only conditioning within block $b$ is on $S_b,I_b$, and its actual
projection $G_b$; unused hypothetical continuations are not conditioned on.
\end{theorem}
\begin{proof}

Let \(S\) contain the first-scale full chart parameters and sweep paths. The preceding construction and
ratios leave \(U_b^*\) product given \(S\), with separate kernels just at \(S_b\). The negligible-scale
nuisance variables appearing in the ratios can equivalently be integrated out. The actual nuisance
continuation trees in uncharted entries, conditional on \(S\), add no dependence on the \(U_b^*\)'s: given
lattice axis paths and sizes these are fresh nuisance fillings; any fixed system of compact ordinary
observations with volumes there converges to its disk kernels for converging positive perimeters and actual
compatible transmitted ranges. Indeed positive macroscopic ordinary steps have \(l/k\to 1\); use
\eqref{eq:original-9}, \eqref{eq:original-12} and compatible path/volume convergence in each fresh disk.
When continuing further given an initial eligible interval, a countable strictly internal target test
suffices, negative loss thresholds avoid exact ends/minima by jump counting, and the child kernels iterate.
This proves the claimed conditional convergence on actual supported first-scale data by finite matched tests
and exhaustion. Ancestor volumes not themselves sampled fresh at this stage are functions by area additivity
in the true limit. Thus microscopic perturbations of the axis parameters from other blocks do not condition
these nuisance macroscopic trees.

Local actual projections, by contrast, can of course condition \(U_b^*\);
they are exactly the separate \(g_b(U_b^*,I_b)\) given \(S\), by the comparison just described. It follows
by disintegration that given \(({\cal T},\text{chart})\), the augmented retained outputs in the
unconditioned chart use the product kernels
\eqref{eq:original-15},
where \(k_b\) is the reference disintegration given \(S_b\), additionally conditioned on the indicated
projection with parameter \(I_b\). Entry modes, color tags, matrix lattice types etc. in this formula are as
specified for the local reference type. One can disintegrate measurably also in the interval parameter:
after fixing the kernel given \(S_b\), use countable refining Borel partitions of the projection range with
ordinary conditional differentiation for each pushed-forward law, defining fallback versions on its null
exceptional data. This works for \(S\)-determined interval parameters without an assertion of continuity of
\eqref{eq:original-15}.

Apply Lemma~\ref{kern:conditional-blocks} to $(S_b,U_b^*)$. The completed
chart parameters give $S_0$, and the actual nuisance continuation trees have
the required kernels depending only on $S$. Lemma~\ref{kern:whole-tree}
identifies the sigma-field generated by these data and the separate $G_b$
with that generated by $(\mathcal T,\text{chart})$. This yields
\eqref{eq:original-15} without continuity of its disintegrations.
\end{proof}

The primary tree $\mathcal T$ is abstract: it contains no ambient conformal
locations. Write $\mathcal E$ for the remaining continuum fields, embedding,
and compatible spatial placements, sampled from their actual conditional law
given $\mathcal T$, independently of all extracted payload and primary
readout arrays and fresh chart marks given $\mathcal T$. We call
\[
 \mathcal X=(\mathcal T,\mathcal E)
\]
the \emph{embedded input}. In spatial assertions below, the full input means
$\mathcal X$. Formula~\eqref{eq:original-15} is unchanged upon this conditional
extension, with conditioning now on $(\mathcal X,\text{chart})$.

\begin{lemma}[Neutral auxiliary charts]\label{kern:neutral-charts}
Fresh smooth chart marks and bounded jump-count selections have their prescribed limiting law
conditional on \(\mathcal T\), jointly with primary readouts. A successful chart selected using only the embedded input and these auxiliary
marks is conditionally independent of the primary readouts given $\mathcal X$.
Consequently its product rule for spatial profiles also holds given $\mathcal X$.
\end{lemma}
\begin{proof}

For clarity, chart choices themselves can be included without altering a primary limiting law conditionally
on its input. Use countable finite patterns with fresh cuts/extra marks sampled with positive smooth
densities, and specified jumps by ordered counting on compact surviving portions. In the original extraction
these marked data have exactly their auxiliary limiting selection law given \({\cal T}\), even jointly with
primary readouts: normalized fraction and time marks are fresh; path prefixes to new ordinary marks use
open-prefix agreement with the primary tests and \(J_1\) jump matching; volume/perimeter labels and
positions of cuts compare as above. Localize counts/cutoffs at continuity tests with bounded multiplicity.
Counting tilts and these restrictions are by input/mark data only. Thus \eqref{eq:original-15} can be
applied on chart successes chosen by limiting data from this library (or absolutely continuous mark choices
conditional on that data). Extra continuum placements and field data are sampled by the actual conditional
law given \({\cal T}\), independently of additional readouts then, and so do not reweight
\eqref{eq:original-15}. This sampling can be done after chart extraction as well; the extra chart cuts/tests
are inputs as just described, and the same joint primary law results.

Explicitly, let $R$ be the primary readout array and let $\Xi$ contain the
neutral chart library and its selection data. Write $C=C(\mathcal X,\Xi)$
for the selected chart. If $V$ is any spatial profile array, a measurable
function of $(R,\mathcal E)$, neutrality gives
\[
 \law(V\mid\mathcal X,\Xi)
 =\law(V\mid\mathcal X,C)
 =\law(V\mid\mathcal X).
\]
On successful charts the middle law is the product law already proved.
The joint law and its marginals are therefore unchanged by removing the
chart, so the product assertion holds given $\mathcal X$. This uses chart
neutrality; an arbitrary mixture of product measures would not suffice.

\end{proof}

\begin{proposition}[Fixed sphere volume]\label{kern:fixed-volume}
The product rule \eqref{eq:original-15}, with the same reference versions at matched local
inputs, holds given the full unit-area sphere input and chart, including after removing Palm
weights. A different fresh nuisance reservoir may be used for each finite-cylinder comparison.
\end{proposition}
\begin{proof}

The product rule persists \textbf{at fixed sphere volume}, with the same reference and kernel versions in
matched local inputs. On such charts leave at least one positive-size unexamined nuisance terminal (a jump
interior will do).

For every cylinder comparison of \(({\cal T},\text{chart})\) we can, \emph{within} it,
further select a fresh jump-filling reservoir unseen by the comparison in its interior. This is the finite-observation exhaustion of Proposition~\ref{cut:fixed-volume}: after any finite compatible paths and
partially observed depths there, continue if needed to select another matched macroscopic birth by strict
continuity tests. Existence and interior compact cutoffs exhaust the unit-area Palm data. Local retained
lattice payloads are never tested in this nuisance reservoir. Integrating it out multiplies the
unconditioned base limit by \(d_u(1-A_{\rm rest})\) by \eqref{eq:original-10} on the corresponding
successes, substituting that volume into ancestor volume entries; \(u\) is the matched perimeter in
\(P=B_n\) units. Joint convergence of the rest volume and labels before substitution is the unconditioned
area comparison above, applicable jointly with the chart limits and the extra birth selection.

To see why the reference versions remain unchanged, let $B$ collect the
macroscopic input outside the reservoir interior, including the local actual
projections $G_b$ and their volume labels. Put
\[
 r(B)=d_u(1-A_{\rm rest}),\qquad
 k_b f_b=\int f_b(v)\,k_b(dv\mid S_b,I_b,G_b).
\]
Let $\tau(B)$ substitute $1-A_{\rm rest}$ into the observed reservoir and
ancestor-volume entries. The unconditioned product rule gives
\[
 \E_{\rm pair}\!\left[r(B)\Phi(\tau(B))\prod_b f_b(U_b^*)\right]
 =\E_{\rm pair}\!\left[r(B)\Phi(\tau(B))\prod_b k_bf_b\right].
\]
Here $A_{\rm rest}$ is measurable from the macroscopic conditioning data:
local block volumes are included in $G_b$, and ancestor volumes are their
area-additive sums. The reservoir is a nuisance filling, so neither $r$ nor
$\tau$ changes $(S_b,I_b,G_b)$.

The reservoir local limit theorem identifies the left side with the
fixed-volume joint limit on the chosen cylinder. The input comparison
identifies the right side with the unit-area input law using the same
reference kernels. One first uses continuous determining probes of the
payload arrays; equality of the resulting finite measures then gives the
identity for Borel probes. This is not a limit through discontinuous kernel
evaluations. The perimeter integration and area change of variables in
Proposition~\ref{cut:fixed-volume} identify the prescribed unit-area Palm
version.
Different cylinder tests can use deeper reservoirs
in the same nuisance piece, and success tests partition/exhaust; there is no need for a conditioned shape
convergence inside a reservoir at an exact specified volume. Convergence of the actual marked input without
passage readouts was already proved by \eqref{eq:original-3}, \eqref{eq:original-10} and exhaustion, also
for these auxiliary chart patterns. We conclude by cylinder determination that given the \textbf{full} unit
input and chart the product rule is \eqref{eq:original-15}. Dividing out Palm bias and restrictions by
limiting input weights as before does not alter it.

\end{proof}

\begin{lemma}[Spatial readouts]\label{kern:spatial-readouts}
The kernels in \eqref{eq:original-15} determine the buffered spatial passage profiles on
protected patches, including universal endpoint and polyhedral extremal-length tests. These
profiles are conditionally independent across separated patches given the embedded
input $\mathcal X$, with
arbitrarily many countable compatible tests in each patch.
\end{lemma}
\begin{proof}

Here and below chart payload formulas are used to localize spatial observations as follows. Inside the
retained pieces one can refine the compatible parts of the augmentation to arbitrarily small meshes (also
partition a sweep further in time, its large attached pieces, and retained rings in length). Use only routes
in the actual projection tree for this purpose. Their labels on the raw maps agree by transmission/open
prefixes with primary comparisons. Thus in a common coupled realization they have the stated upper images at
continuum locations, including their incidence stars; compact constraints pass into the covering pieces and
those meeting a prescribed compact set have all images in any strict open enlargement eventually. Together
with the exact local gluings this lets either the primary cover arrays or the localized arrays evaluate the
\emph{same} spatial regularizations.

More explicitly keep attainability and exclusion with compact
ports/ranges versus open enlargements (and similarly ordered/subpath marks and budget slack); between any
two strict margins a finite cover by either refining alphabet gives intermediate tests for all paths. For
universal endpoint tests apply both containments also to all lifts/vertices, with the corresponding opposite
margins for allowed connections. Curve-family costs such as polyhedral extremal lengths with compact versus
open passage constraints are squeezed by inclusion, using the true oriented flags near the confinement;
finite ordered tubes/ports can enforce robust crossings and winding. A family confined there uses no
exterior conformal metric densities in the extremal-length supremum. Taking one-sided envelopes over margins
or testing at continuity contours gives determinate readings from either array (subsequence upper versus
strict-lower comparisons likewise use the corresponding margins). This explains why exact boundary-equality
readouts in an exterior alphabet are not needed.

Arbitrarily many countable such tests strictly in an inner
protection can be treated together. Their conditional laws use \eqref{eq:original-15} and the actual
protected positions of those compatible refined pieces, and in particular are independent across patches
given $\mathcal X$, by the neutral-chart identity. Such array extractions can always compactify costs or use countable hyperspace tests;
path tests need not yet have a vanishing endpoint-cost modulus.

\end{proof}

\begin{lemma}[Conversion of fixed units]\label{kern:unit-conversion}
The same lattice reference experiments compare scales in any fixed positive limiting ratio.
For \(P\mapsto\lambda P\), time units change by \(m_{\lambda P}/m_P\to\lambda\) and area
units by \(\lambda^2\). With common or explicitly converted cost units, the local kernel
comparisons require no scale invariance of the passage law.
\end{lemma}
\begin{proof}

These comparisons of versions are intrinsic to the same lattice experiments. They allow replacing the
perimeter scales by others in fixed positive limiting ratio with the macro data converted accordingly:
\(m_{\lambda P}/m_P\to\lambda \), and the triangle count units (using asymptotic inverses of \(B_n\))
transform by \(\lambda ^2\). Reference experiments with those normalizations are smooth lattice changes of
density of the same experiments after converting the coordinates, with ratios depending only on the
conditioned parameters on compact ranges; raw words, runs and terminal laws at exact sizes are identical.
Readout cost units of course have to be used commonly or separately converted. Thus constant local
field/coordinate comparisons below involving fixed unit conversions do not assume passage-law scale
invariance.

\end{proof}

\subsection{Matching charts over exterior extensions}

At bounded local complexity there are finitely many integer lattice and residue
types, and each fixed chart has its reference conditional kernels. We still need
common finite local complexity over the exterior extensions. We must also match
the \emph{values} of the conditioned macroscopic parameters: a finite list of
matrix types alone would not give finitely many kernels. The following chart
construction supplies both conclusions. Work with a compact inner test and a strictly larger open field buffer, in the previous
common local wall comparison of Proposition~\ref{wall:finite-local-patterns}, with the two designated ends outside (and central loop separating them). Use
countably many conditional extensions, sharing the local fields including actual heights. One compares a
\emph{marginal} kernel this way; \eqref{eq:original-15} for simultaneous separated observations only needed
a successful chart for the one overall drawing. Several fixed perimeter-unit conversions and countably many
common cost extractions can all use the same order comparison of Proposition~\ref{wall:simultaneous-targets}.

\begin{lemma}[Private source pools]\label{kern:source-pools}
For countably many conditional exterior extensions of common local fields, one can couple fresh
source choices by private local pools. Averaging the pools gives the correct conditionally
independent source law in each full input. For every fixed positive diameter threshold, the
restricted caller keys needed by the coupling have finite conditional support.
\end{lemma}
\begin{proof}

Use extra independent local random pools for the source positions. These can be chosen as functions of local
data and iid uniform seeds. For detail, use a still buffered region inside field agreement, large enough for
all local observations below. A caller whose boundary meets it has a key consisting of the caller side/open
domain and boundary there (not its remote portions), and the side length measure there. Conditional key
types for boundaries reaching that region can be enumerated measurably as atoms. Indeed for callers of
diameter at least any fixed positive bound there are only finitely many such restricted keys over countably
many conditional copies, by the depth/cell comparison in a finite cover: use small balls relative to the
bound with enlargement, so that a boundary there is a local operation in a caller extending across the small
ball's comparison collars. Away from boundary the open-domain bit only has local whole/empty choices. Group
the restricted pieces by finite tags. Arc lengths agree by the boundary comparison. Thus conditional iid
sampling already gives finite support at each diameter threshold; enumerate atoms, for example using a
countable weighted sum law of keys at that threshold. Smaller whole internal hierarchies can equally be
matched by their own domains.

Two distinct caller keys within a copy (with boundary strictly here) need not
share a uniform seed: the open-domain portions distinguish disjoint interiors, and strictly nested callers
have disjoint rims. Use a separate key for a central-ring phase if it uses a position there.

Couple by
choosing whether a length source lies on the observed boundary portion by its actual conditional probability
and using the pool point sampled on that portion if so, otherwise using an exterior point. Caller geometry
uses the unrouted canonical hierarchy, not earlier source choices. Thus after averaging the pools the
sources are correctly conditionally independent. This works also for many disjoint patches with private
pools, choosing among disjoint boundary categories over those patches by the actual full-field
probabilities. Extra pools can be conditionally neutral for passage draws given the full input and actual
sources. For the following comparisons we can add dense local target pools too, with private completions;
positive joint length-density marginal choices in any finite chart suffice.

\end{proof}

\begin{proposition}[Common finite local chart data]\label{kern:matching-charts}
Fix a compact inner test with a strict field buffer. In the conditional exterior comparison with
private source pools, successful mosaics can be chosen so that the retained raw-piece instructions,
local open geometry and order, interface parametrizations, actual entry projections, and effective
macroscopic gluing shifts range over one finite collection. The collection includes all closed
viable matching cases and does not depend on the cost units or the readout extraction. This
collection is finite almost surely for the common local data and pools; no deterministic bound
on its size is asserted.
\end{proposition}
\begin{proof}\leavevmode

\paragraph{Expanding callers and crossing walls.}

We claim that for buffered inner readings the vector of conditional laws (given full inputs) as evaluated
over these extensions has only finitely many possibilities, with these pools, almost surely. Use successive
nested buffers \(K\Subset D_1\Subset \cdots \Subset D_{10}\Subset U\) about the protected compact \(K\)
(extra intermediate margins implicit), inside the region of pool comparison where needed. The numbers here
only specify ordered clearances. Expand all canonical callers whose closures meet both \(\overline D_2\) and
\(D_4^c\), reaching them from the two central sides as needed; per sample there are finitely many by the
meshes. Among these, instances with any local operations in the larger buffered charts under use have
bounded local number and finitely many local open shapes/half-layer and boundary order possibilities by Proposition~\ref{wall:finite-local-patterns} (again one can cover by smaller boxes than the required diameter).

Other
enclosing ancestors can be remote-only, with no first-fan or rim operation in the corresponding protected
band. The relative half-layer shifts furnished by different local shields are allowed to differ.
We use their finite local geometric and order patterns separately; Lemmas~\ref{kern:integer-splitting}
and~\ref{kern:finite-lattice-types} account jointly for the resulting integer translations.
No common depth origin for all shielding neighborhoods is required.

In each expanded caller demand births of all its similarly expanding direct caller children at that
stage.

Also when comparing full-path obstacles in the intervening band, one can demand births of eligible
positive terminal walls there that reach out of the larger path-comparison region (e.g. boundaries meeting a
neighborhood of the closed middle band and also reaching outside a slightly larger neighborhood within
\(D_6\)). There are only boundedly many such extra crossing walls in local-relevant calls, uniformly over
the extensions. Expanding child instances with boundaries approaching a test likewise have the finite
classifications. Large disjoint children reaching an inner protection but with their whole rim staying
outside a larger buffer are bounded in number per parent also by interior ball packing.

\paragraph{Demanded births and common target choices.}

Each demanded canonical birth can be made on a single parent-fan branch to an ordinary target in an open
interval. Its pinch/discovery \(p\) is strictly in the original caller \(C\), hence on an interior gap with
no target split along it. Indeed loss strictly into untouched original labels puts the loss threshold and
tip on the true original boundary. Use an exterior target if a positive open portion allowing the birth does
reach the exterior (here exterior targets for traversal comparison may just be beyond \(D_7\)). Otherwise
the open interval for this gap is local. There are finitely many shared pool choices sufficient in that
latter case:
\begin{itemize}
\item For \(p\) within a larger local chart (\(D_8\) with enlargement), a demanded macroscopic canonical child's boundary there stays positively separated from \(\partial C\), uniformly for the matched finite geometries. Complete positive walls obey this too. Thus use interior gap markers there; only finitely many needed, by the common crosscut pattern on interior compact parts and gap finiteness, to have a marker on the very gap in use.
\item If \(p\) is remote while its gap target interval is trapped in the local-target region, the starting endpoint of the excursion (a current extreme contact bounding that interval) is there as well. The gap passage to \(p\) thus crosses an observable collar without a split, giving an interior marker by the traversal result.
\end{itemize}

All paths using the corresponding marker share this gap/prefix; its open target interval on the local
boundary is exactly among the finite options already described. This proves the claim about pool choices
(enlarge charts for endpoints on transition buffers). Remote-only ancestors instead need just exterior
positions and no shared local gap tests where operations are absent.

\paragraph{Eliminating passing unstabilized components.}

Additional paths needed for fineness in a caller can still use finitely many local types. Tag the finite
full-fan/traversal alternatives, including feasible exterior-target sequences of buffered pieces. In a
representative choose sufficiently many ordinary full target paths so that no untouched \emph{unstabilized}
negative goes across the intervening band, by the finite-target stabilization. Require their local
portions/obstacles there also in actual copies of this case. Local targets can be chosen from common pools.
For exterior targets use counterparts with the corresponding buffered anchor-piece sequences; it suffices to
require a prefix through the needed pieces, which for a tested exterior-terminating path visiting the
smaller band occurs for positive-length choices (last passages there strictly before terminal). Feasibility
of such exterior choices can itself be tagged; actual further pieces are covered by the same finite
traversal classification.

Distinct desired paths/constraints can use positive joint density marks, not exact
coincidences between fresh marks. Positive obstacles wholly local that are hit along the compared passages
are drawn there (or already discovered); ones affecting the band but first hit farther out are included by
the extra crossing-wall demands above.

Thus an open surviving uncharted negative traversing from \(D_1\) out
of \(D_5\) after the full chosen paths, in the same caller, would have an interior subpassage from \(D_2\)
to outside \(D_4\) avoiding all representative obstacles on that band. In the representative such a
complementary passage is inside a stabilized canonical first-stage child. By the shared local domain and
half-layer patterns it is then in such a child's interior in the actual drawing there too (the subpassage
can use strict intermediate margins). This child is expanding and targeted, contradiction to an untouched
unstabilized negative. Expanded children themselves are assigned further calls.

In a portion wholly covered
without any local first-fan operations a crossing subpassage similarly has to lie inside a targeted crossing
child; no bound on the number of remote generations is used. This gives the requisite exclusion. All the
representative demands refer to finite matched local-pattern cases for expanding parents with operations
there; remote-only parents need no copied local target path. Additional suppression of very large leftovers
away from the comparison is unnecessary.

\paragraph{Truncation after forks and local terminal traps.}

Truncate the finite families before terminal after needed discoveries and forks, leaving only nonpassing new
terminal polygons. At an exterior target cut sufficiently late off the inner comparison; at a local target
the final segment is in matching charts. It can be trapped near that target by a complete small crosscut of
this segment: boundary contacts on both original sides accumulate as proved above, so there is a late
excursion with endpoints on opposite sides, cutting off the small Jordan-coordinate shadow toward the
target. This trap works throughout the corresponding matched terminal-segment geometry.

For local targets
whose end regions will be tested choose it avoiding other distinct target positions (only finitely many
shared choices, exterior targets away) and small enough to be past macroscopic demands, then one can match
sufficiently late clock cuts. If an end and late forks are all away from the region needed for retained
pieces, their later truncations can instead be private there. Forks between the finite paths are ordinary
preterminal designated splits; preserve common prefixes and include these in the finite operations.
Truncating \emph{after} them only collapses tail subdivisions into the indicated target polygons.

\paragraph{Significant jumps and finite groupings.}

We give more subdivision details to clarify what information can be matched. Along the finite chosen paths
designate individually all jumps of significant diameter near the local observation (for example all with
closure meeting \(\overline D_4\) and diameter above a small fixed tolerance relative to buffers), in
addition to needed forks/births. Use their two chart branches, terminal assignments or new callers as
scheduled. Uncontinued crossing disks near \(D_1\) reaching outside \(D_5\) have been excluded by the
preceding target choice.

There are uniformly only finitely many local significant-jump portions/groupings
for these finite traversal types. Indeed the local obstacles of a full chosen path are its trunk, initial
boundary and hit positive walls. Which crossing positive walls have been hit or were already hit earlier as
needed on buffered anchored pieces needs only finite tags (terminal portions for shared local targets are
themselves matched). Smaller ones reaching the observation have visible contacts in the larger charts.
Negative jumps are actual complementary components, with pinch the last trunk contact, by the exact loss
rule; positives have first contact at discovery. Large components' local pieces going across collars admit
only finite exterior groupings when these local obstacles match.

One way to check finiteness is in a
representative: global branches have finitely many components above each diameter and continuous closed
component maps; in any one such domain interior points in a compactly inset local region cannot belong to
infinitely many different components of its restriction to the larger region (join convergent preimages
inside the prime disk using continuity). Small components wholly inside agreement regions have no exterior
regrouping. Use also balls smaller than the diameter threshold in a finite cover.

A large jump's \(p\), if
local where needed, is thus identified among finitely many positions on the matching ordered path pieces, by
first/last contact on the corresponding boundary. Similarly a local fork of two selected paths is identified
by the common-prefix and no-reconnection rule. Order, sharing and accessible side for the resulting finite
local operations can all be tagged.

\paragraph{Retained sweeps, rings, and distinguished slots.}

Between designated operations take additional ordinary cuts, sufficiently fine on relevant compact portions
of the paths, so that a sweep meeting the inner protection is itself retained whole in a still buffered
matching neighborhood, with its nondesignated attachments. This uses continuity of trunks and the small
diameters of all those other jumps near the buffers. It requires only boundedly many local operation cuts,
whose positions can be matched: portions visiting the smaller bands are in the finitely matching ordered
path segments described by the anchors or local source/target segments above. Cover their needed time
portions with room inside agreement, subdivide about local designated times there, and make cuts within that
coverage before reaching its outer edges. Whole intervals retained then have both ends at matched positions.

Extra intervals away can have private cuts. Terminal pieces if retained encompass their entire disk in a
matching region: for a jump or an internal caller the whole boundary/domain and root are then known by the
compared path/domain choices; for a cut-time post polygon use the chosen terminal traps and the known
ordered obstacles up to the compared cut, with the tip's access. At ordinary times there is a single access
at the tip by the contact rule of Proposition~\ref{cut:critical-cutting}.

Ring chunks (central included) likewise take small buffered arcs,
putting additional circumference cuts \emph{within} the shared arc portions so that retained chunk sizes
need not count back to an exterior first position. Point attachments and radial joins across local adjacent
chunks are included. All retained pieces touching the needed inner part thus fit in a strict interior buffer
such as \(D_6\). Local sweeps and chunks not touching the path tests can be included for parameter
calculations while staying buffered. Source-index routes in the integer tests can use still smaller
protection (strictly inside \(D_1\), retaining all operations potentially involving those points); arcs back
to their local anchors just need the larger matched buffer.

The number of distinguished slot prescriptions
for the weighted accounting on a further enlargement of these retained supports is bounded as well: fresh
marks and chunk cuts there use only the finite local-reaching instances; actives, paired bridge/birth slots
and their immediate neighbors are at the corresponding tip operations. A persistent local target of many
remote operations is the same slot. Use no gratuitous origin inside this buffer at remote-only steps.

\paragraph{Matching clock cuts with positive density.}

The clock cuts just used can be matched with positive-density marginal choices conditional on the
macroscopic data. Along a common directed portion the clock is read by small positive-jump counts by the
Cauchy law; giant jumps with nonlocal data, individually designated, do not affect elapsed time comparisons
by such counts. Smaller apertures used on the retained portions match with their length measures. Use for
instance common diffuse time cuts with room on shared portions, with arbitrary private subdivisions outside
(translations in elapsed time from a possibly remote start are allowed). Circumference cuts similarly have
ordinary densities. Existence with room in the finite local cases suffices to use such marks; counts/cuts
can be increased on the common segments using a countable randomized dense library. Each individual chart
obtained this way need only be absolutely continuous as a marked choice conditional on full inputs relative
to the auxiliary chart tests, since those tests themselves are conditionally neutral for the primary
limiting readouts. Segment ends at chosen jumps use the jump-count rather than an independent continuous cut
at that exact position.

\paragraph{Sides, increments, and suffix parametrizations.}

On a retained ordinary sweep, all nondesignated positive and negative disk boundaries, order, sides and
lengths in it are now locally matching. Small negatives wholly inside the matching region are the actual
components as just discussed. Side can be read by oriented traversal there, not a new exterior bit at each
such jump. For example a negative jump disk lies on its indicated side of the full simple trunk chord (as
also seen in the side-height quotient); its boundary has an interior trunk contact. Otherwise its connected
Jordan boundary would be covered by pairwise disjoint positive circles and the caller rim alone, hence lie
on just one by the countable closed partition theorem, impossible. For this use only its
circle case is needed: a circle cannot be partitioned into countably many disjoint nonempty
closed proper sets. To verify that case, write a putative partition as \((F_n)\), let
\(U=\bigcup_n\operatorname{int}(F_n)\), and put \(E=\mathbb S^1\setminus U\).
If \(E\) is empty, connectedness of the circle already forces one \(F_n\) to be the whole
circle. Otherwise \(E\) has no isolated point: the two adjacent arcs around such a point
would belong to interiors of partition sets whose closedness forces both sets to contain that
point, making it interior as well. Each \(F_n\cap E\) has empty interior in \(E\). Indeed,
if an arc met \(E\) only in \(F_n\), choose three ordered points \(x<z<y\) of \(E\) in
that arc. Every component of \((x,y)\setminus E\) lies in some \(\operatorname{int}(F_j)\)
and has endpoints in \(F_n\); closedness and disjointness force \(j=n\). Hence
\([x,y]\subset F_n\), contradicting \(z\in E\). The nonempty compact space \(E\) would
therefore be a countable union of closed sets with empty interior, contrary to Baire's theorem.

At an open-caller trunk contact the accessible local side relative to the oriented simple crosscut gives the designation. Positive
circles likewise use orientation/sign and trunk.

The retained increments in time then agree as well by
symmetric Cauchy summation of the jumps on compact surviving portions. Suffix parametrizations from the
active tops are local: a.e. exposed height on input first pops strictly on a negative aperture in that
sweep, whose boundary is matched; output suffixes similarly trace backward to positive insertions, then
extend by continuity. Whole retained terminal parametrizations use the paired side lengths in the known
domain and active/attachment prime coordinate (equivalently strict inherited intervals on original and
positive arcs up to null sets). Thus these interface maps including their threshold points agree.

\paragraph{Local entry projections and protected positions.}

In the resulting uncharted local entries also \(I_b,G_b\) for \eqref{eq:original-15} are local up to the
finite alternatives. The untouched original open interval comprises the cycle labels still on the original
caller rim, up to length null sets: strict backward transmission a.e. to the original rim or an earlier
positive aperture shows this, since inserted positive walls on the branch in that caller are disjoint from
its rim. Take the transmitted open interval in the prescribed cyclic order. Continued original target paths
there are known by the full local-target comparison (common ordinary targets/open prefixes suffice, stopping
the shared incoming part at the matched entry).

Later restart domains within the retained fillings have the
shared internal hierarchy, field lengths and coupled actual sources. True time-cut post domains used whole
can employ the prime boundary of the matched open domain; the field inside or the strict inherited labels
give their length measures. This matches also the protected embedded positions used with refinements, not
just abstract total lengths. Internal subtrees on a compared complete local caller require no exterior
traversal comparison of their own.

\paragraph{Finite values, including singleton cases.}

We highlight one detail about the \emph{values} of the gluing shifts over exteriors. All tested local
thresholds when binding are at finitely prescribed locations for these charts, including a matched sweep
minimum or wrap/activation location. Each such point has only finitely many preimage labels on a fixed
source interface. For a frontier in one path from a canonical Jordan caller, this is
Corollary~\ref{cut:finite-frontier-labels}: a non-tip class has a single direct linkage height on its side, and at a tip there is
only that real time with its boundedly many side copies (possibly a jump fold), with at most one nontrivial
side linkage. Circle endpoints cause just the stated identifications. Current-cycle frontiers after forks
toward an inherited target can use its path from the original caller. Chunks on loops themselves are
Jordan-parametrized.

Thus the binding offsets (counting back to the top anchor if needed) have finite exact
macroscopic possibilities. The source-coordinate closed intervals are consequently cut at only finitely many
possible values; on a final viable matching interval choose one of its known ends or a common interior
value, and again there are only finitely many preimages at the consumer. This fixes the matching translation
up to finite choices, \emph{including singleton cases}. Remote steps with no binding on a viable prefix add
no subdivision or exterior threshold value. This explains why collapsing the potentially long route list to
the effective local tests does not insert an extra free phase.

\end{proof}

\subsection{Finite vectors of kernels and changes of field class}

\begin{theorem}[One finite alternative for a countable kernel vector]\label{kern:finite-vector}
Let \(\mathcal A\) be any countable family of common cost extractions and fixed unit conversions.
Conditional on the local real fields (including their heights), the local auxiliary field, and
the private source pools, the vector of full-input conditional profile laws
\((K_\alpha)_{\alpha\in\mathcal A}\) has finite support. Equivalently, one tag \(T\) with finite conditional support almost surely suffices to write every coordinate as \(K_\alpha=F_\alpha(L,T)\), where
\(L\) denotes these common local data. The tag is common to all coordinates; its number of
possible values may depend on \(L\).
\end{theorem}
\begin{proof}

Combining \eqref{eq:original-14}–\eqref{eq:original-15} and this construction, the local reference readout
instructions, lattice and residue types, conditioned values and embedded comparison alphabets, as needed to
compute the robust profile kernel, have only finitely many possibilities over the copies. For this
bounded-alternatives comparison one may work in arbitrary-area pair law (no specified nuisance volume
reservoir is needed). Omitted remote histories can change the matrix but the kernel uses only its local
projected type as proved above. All countable common unit/cost conversions use the same finite
geometry/order choice. Hence conditional on the local fields and independent source pools the corresponding
vector of conditional laws is finitely supported. This conditional-support argument can use conditional iid
exterior samples and the matching extra target/cut libraries just described: all marginal
\eqref{eq:original-15} identities still hold, and the original conditional law given full input does not
depend on those extra library draws. Finite support for the marginal vector then follows already by
finiteness over the countably many iid extensions. Feasibility can always mean positive conditional weight
under the particular common reference measure class.

To identify the common chart choices, record the geometry and order choice of
Proposition~\ref{kern:matching-charts}, the effective local reference type of
Lemma~\ref{kern:finite-lattice-types}, and the conditioned local macroscopic values from that
proposition. Whole retained words and fillings are the reference payloads; every refining
readout for every \(\alpha\) is extracted from those same payloads. Duration and microstep
counts belong to the sampled parameters and raw words, not to a new extraction-dependent
constant term in \eqref{eq:original-14}. All closed viable cases have already been retained in
Proposition~\ref{kern:local-projections}; their exact microscopic pass/fail results are payload
readouts. Neither changing a budget nor changing units adds a new geometric choice. Thus, after
a single countable diagonal extraction of the references, \(F_\alpha\) may vary with \(\alpha\)
but the argument encoding these choices ranges over the same finite set in the coupled
comparison. The conditional-support conclusion below then permits \(T\) to be chosen simply
as a measurable index of the atoms of the vector law given \(L\).

For completeness, the conditional iid argument proves support finiteness, not just finiteness
within a selected list of examples. A probability measure on a standard Borel space that is not
supported on finitely many points gives an iid sequence infinitely many distinct values almost
surely. This follows either by repeatedly sampling outside any finite set of previously seen
atoms, or from its non-atomic part. Apply this fact conditionally to the countable product of the
kernel spaces. The matched charts above bound the distinct vector values among all the
conditional iid exterior copies by a finite number. Hence the conditional vector law itself
has finite support.

\end{proof}

\begin{corollary}[Uniform fixed-label candidate tests]\label{kern:candidate-draws}
In a fixed prototype law, the number of candidate kernel vectors and the reciprocal of their
smallest positive conditional mass can be truncated once in probability. Arbitrarily high
marginal success at a fixed label therefore gives arbitrarily high success for draws from all
candidates, with a loss bound uniform in the label. The conclusion concerns each fixed label;
it does not require simultaneous success over infinitely many labels. Across separated patches,
independent arrays of candidate draws can include the actual selected draws.
\end{corollary}
\begin{proof}

One consequence for high-probability uses is that the number of candidate kernels and their smallest
positive conditional weight can be truncated in probability in a fixed prototype law,
\textbf{simultaneously} over the labels in the vector. Thus high probability with arbitrary slack can be
required for every candidate's draw, at a fixed label (uniform loss bound across labels), using the local
fields, pools and auxiliary sampling. Under absolutely continuous transport the actual choice stays among
the candidates. Across disjoint patches with pools as above one can have independent arrays of candidate
draws conditionally on input and pools (at respective tested labels), whose distribution then needs just the
separate local fields/pools, with actual profile draws occupying the selected slots by the product rule.
Exterior variables can select alternatives but are not used in the simultaneous all-candidates success test.
These are statements at the limiting kernel level.

Indeed let the conditional vector law have atoms \(\kappa_1,\ldots,\kappa_N\) and positive
weights \(p_1,\ldots,p_N\), and put \(p_* =\min_j p_j>0\). For every \(\eta>0\), choose
\(N_0<\infty\) and \(\delta>0\), using this one vector law, so that
\[
 \Prob\{N>N_0\text{ or }p_*<\delta\}<\eta.
\]
If \(q_{j,\alpha}\) is candidate \(j\)'s conditional failure probability at label \(\alpha\),
then on the complementary event
\[
 \sum_{j=1}^{N}q_{j,\alpha}
 \le \delta^{-1}\sum_{j=1}^{N}p_jq_{j,\alpha}.
\]
The union bound gives an unconditional all-candidate failure probability at most
\(\eta+\delta^{-1}q_\alpha\), where \(q_\alpha\) is the marginal prototype failure
probability. The same \(N_0,\delta\) work for all labels, although the success event is tested
at one label at a time. Absolute continuity transports the supported alternatives; density
truncation transports the probability estimate when needed. For separated patches, first draw
the arrays independently given their private local data and pools, then identify the externally
selected slots. The product rule supplies the joint law of those selected entries.

\end{proof}

\begin{lemma}[Exact comparison between field classes]\label{kern:field-comparison}
In supported arbitrary-area field classes, common auxiliary geometry and matching real fields
on observation neighborhoods, or exact constant differences there with converted units, give
identical local profile kernels on the same successful finite charts. A bounded smooth Weyl
change outside these neighborhoods and absolutely continuous choices of the finitely many
relevant exterior sources preserve the comparison.
\end{lemma}
\begin{proof}

Here also are exact comparison conventions between field classes. Randomize area and embedding (e.g.
cylinder translation and rotations preserving two marked ends, and further independent Möbius changes with
positive densities if needed). By the Palm split already proved, a positive-density mixture of perimeters
and pair bases can be transported to the arbitrary-total-area sphere class with a separating central loop
under positive counting weights; give every feasible separating-loop choice positive support when comparing
classes. Color parity, length phases and internal source data use the conditional rules above. Bias/unbias
by full limiting inputs is harmless for \eqref{eq:original-15}.

If two full inputs have common \(J\), ends
and loop choices and their real fields match on observation neighborhoods (or differ by exact constants
there with matching unit conversions), differing by a bounded smooth Weyl change elsewhere, successful
\textbf{same} finite unclocked charts give exact comparison of kernels on those observations. This assertion
is for coupled typical marginal inputs in the supported classes: caller sources can be shared with
absolutely continuous marginal choices. Indeed only finitely many noninternal callers influencing a strictly
protected chart have to be serviced (all required ancestors included); all further sources relevant inside
retained blocks are in whole constant-comparison neighborhoods. Unrelated changed subtrees need not be
coupled.

Boundary lengths multiply by the continuous boundary factor, by the disk-field/arc comparisons
above. Clock cut measures on common unclocked paths are equivalent by the positive-jump count formula and
uniform upper/lower bounds on that factor; on retained segments elapsed clocks compare exactly in converted
units. Thus the needed finite outer choices have equivalent supports with positive densities given the
fields; inside subtrees one can share the source law exactly.

The very same affine histories/thresholds can
be used: unchanged-slot viable transfers compare by their interval placements along the actual cut path,
which do not change under remeasuring arcs, and near protected points their required shift values and
binding arc lengths are the same in matched units. More explicitly the interval arrows carry actual arc
labels together with their closure endpoints; remeasuring by a continuous positive position weight maps
these monotone parametrizations homeomorphically onto the new ones, consistently on unchanged intervals and
at every specified prime-end threshold. Cut clipping can shrink with room in the common local coordinates.
Protected \(I_b,G_b\) and local subdivisions remain matched. A whole finite chart may be shared this way
(include a nuisance reservoir when needed), so no assumption of uniqueness of an exterior-order choice from
local data is required for this comparison.

\end{proof}

\begin{lemma}[Local transfer to area one]\label{kern:area-one-transfer}
Local almost-sure kernel assertions in the arbitrary-area sphere class transfer to the unit-area
class on compact slabs away from the marked cylinder ends. The comparison preserves the real
field on the slab, uses the fixed-volume product rule, and remains valid after restoring area
marks and the three-point normalization.
\end{lemma}
\begin{proof}

For local almost sure transfers to the area-one version work before three-point normalization, off the
marked cylinder ends, in the representation \eqref{eq:original-2} with a random translation. There is a
useful exact way to match to inputs absolutely continuous in the arbitrary-area class while preserving the
real field on any compact interior slab. Choose a nonzero smooth Cameron–Martin function \(g\) of the
lateral field with zero angular mean, supported away from that slab. Write the lateral Gaussian as an
independent Gaussian scalar \(u\) times \(g\), plus its complementary part. Conditional on the latter and on
average-process data, the normalizing area \(A(u)\) before constant subtraction is strictly positive
analytic and strictly convex as a function of \(u\): it integrates \(e^{2u g}\) against the other positive
finite area by continuous chaos correction, with full support. Hence \(-(\log A(u))/2\) has a density.

Resample the coefficient of \(g\) independently but \textbf{keep} this old normalizing constant. The output
equals the old field throughout the slab and has full-field law absolutely continuous in the arbitrary-area
class (conditional on the complementary variables its coefficient and constant now have a joint density).
Keep \(J\) and separating-loop geometry, using the positive-choice and source comparisons just explained.

This transports supported local kernel assertions to area one by the fixed-volume rule
\eqref{eq:original-15} on its own charts and the exact matching comparison; one does not condition a generic
full-input almost-sure metric assertion on a null area event. Similar local comparisons on separated inner
tests can use wide encompassing slabs. The sampled embedding marks can finally be restored using area
marking.

The density assertion used here follows directly from one-variable calculus. Conditional on the
complementary field, write
\[
 A(u)=\int e^{2u g(z)}\,\mu_0(\mathrm dz).
\]
The bounded smooth function \(g\) and finite measure \(\mu_0\) make this function real analytic
on \(\mathbb R\), by termwise differentiation on every compact \(u\)-interval. Moreover
\(A''(u)=4\int g(z)^2e^{2u g(z)}\,\mu_0(\mathrm dz)>0\), since \(g\) is nonzero and
\(\mu_0\) has full support. Thus \(A'\) vanishes at at most one point. The map
\(u\mapsto-\frac12\log A(u)\) has nonzero derivative off that point and is locally a smooth
change of variable there. A nondegenerate Gaussian coefficient consequently pushes forward
to an absolutely continuous law for the retained constant. The resampled coefficient is
independent of this old constant conditionally on the complementary variables, so their pair
has a joint density. This is the precise density mechanism; no conditioning of an
arbitrary-area null event is involved.

\end{proof}
